\section{Axiom 路线的关键判断：为什么要做研究型公司}

Axiom 这样的 Neo Lab 不是普通应用公司。它要面对一个很深的产品悖论：数学研究的最终用户很少，短期商业化不如通用 AI 应用清晰；但一旦系统真正提高数学研究速度，它可能影响科学、工程、密码学、物理、材料和 AI 自身。高风险和高上限同时存在。

\subsection{Ken Ono 事件的信号}

媒体叙事喜欢讲“57 岁终身教授给 24 岁创始人打工”，但更值得关注的是信号本身：顶级数学家愿意把时间放到 AI for Math 上，说明他们认为这不是玩具。数学家加入带来的不是宣传价值，而是问题选择、研究审美、证明质量判断和社区连接。

\begin{importantbox}{为什么数学家不可替代}
AI 可以生成证明候选，但数学家决定哪些问题重要、哪些证明有解释力、哪些方向值得投入。AI for Math 若脱离数学家共同体，很容易退化成题库自动化。
\end{importantbox}

\subsection{融资不是结果，而是 runway}

Axiom 的大额融资说明资本愿意押注 AI for Math，但融资只是 runway。真正要验证的是：它能否构建可持续的定理库贡献、证明自动化工具、数学家工作流和研究协作产品。否则融资规模只会放大压力。

\begin{warningbox}{Neo Lab 的风险}
研究型公司容易两头不靠：学术界认为它太商业，商业世界认为它太基础。要穿过这个夹缝，需要一个能被验证的中间产物，例如形式化证明效率提升、定理库增长、数学家协作平台或高价值研究突破。
\end{warningbox}

\subsection{本章小结}

Axiom 的路线不是“做一个数学聊天机器人”，而是试图把数学家、形式化系统和模型训练组织成研究机器。它的风险和价值都来自这个定位。

